\documentclass[12pt]{ctexart}
\usepackage{amsmath, amsthm, amssymb, bm}
\usepackage[indent=0pt]{parskip}

\newcounter{probid}
\newcounter{subprobid}[probid]

\author{杨濛 2021213247}
\title{数学作业\\第一周}
\date{\today}

\newenvironment{problem}{
    \stepcounter{probid}
    \noindent\textbf{\arabic{probid}. }
}{\par}
\newenvironment{subproblem}{
    \stepcounter{subprobid}
    (\arabic{subprobid})
}{\par}
\newenvironment{solution}{
    \noindent\textbf{解: }
}{\par}

\begin{document}
\maketitle
\begin{problem}
    写出下列级数的通项,并表示成\(\sum\limits^\infty_{n=1}u_n\)的形式:
\end{problem}
\begin{subproblem}
    $\dfrac{2}{1}-\dfrac{3}{2}+\dfrac{4}{3}+\cdots
    = \sum\limits^\infty_{n=1}\left(-1\right)^n \dfrac{n+1}{n}$
\end{subproblem}
\begin{subproblem}
    $x - \dfrac{x^3}{3!} + \dfrac{x^5}{5!} +\cdots
    = \sum\limits^\infty_{n=1} (-1)^n-1 \dfrac{x^{2n-1}}{(2n-1)!}$
\end{subproblem}
\begin{problem}
    用级数敛散定义判定下列级数的敛散性.如果收敛,试求出级数之和.
\end{problem}
\begin{subproblem}
    $\sum\limits^\infty_{n=1} \ln\dfrac{n+1}{n}$
\end{subproblem}
\begin{solution}
    \[
        \sum\limits^\infty_{n=1} \ln\dfrac{n+1}{n} 
        = \sum\limits^\infty_{n=1} \left( \ln (n+1) - \ln n \right)
        = \lim\limits_{n\to\infty} ln(n+1)
    \]
    显然这是一个无穷极限,故原级数发散.
\end{solution}
\begin{subproblem}
    $\sum\limits^\infty_{n=1} \dfrac{1}{n(n-2)}$
\end{subproblem}
\begin{solution}
    这个级数的部分和为
    \begin{equation*}
        S(n) = \dfrac{1}{2} \left( 1 + \dfrac{1}{2} 
        - \dfrac{1}{n+1} - \dfrac{1}{n+2} \right)
    \end{equation*}
    此时有$\lim\limits_{n\to\infty} S(n) = \dfrac 3 2$, 故级数收敛至$\dfrac 3 2$.
\end{solution}
\begin{subproblem}
    $\sum\limits_{n=1}^\infty \sin\dfrac{n\pi}{6}$
\end{subproblem}
\begin{solution}
    显然对于任意一个$n\in\mathbf{N_+}$,$\sin\dfrac{n\pi}{6}$
    与$\sin\left(2\pi-\dfrac{n\pi}{6}\right) 
    = \sin\dfrac{(12-n)\pi}{6} 
    = -\sin\dfrac{n\pi}{6}$均是级数的项,所以
    \[\sum\limits_{n=1}^12 \sin\dfrac{n\pi}{6} = 0\]
    于是易知其部分和
    $S(n)\sum\limits_{n=1}^n \sin\dfrac{n\pi}{6}$
    是一个周期函数,其极限$\lim\limits_{n\to\infty}S(n)$发散.
\end{solution}
\begin{problem}
    判断下列级数的敛散性:
\end{problem}
\begin{subproblem}
    \(\sum\limits^\infty_{n = 1} \left(\dfrac 1 2\right)^n\)
\end{subproblem}
\begin{solution}
    这是一个等比级数, 其部分和为
    \begin{equation*}
        S(n) = \sum\limits^\infty_{n = 1} \left(\dfrac 1 2\right)^n
        = \dfrac{1-\left(-\dfrac{1}{2}\right)^n}{1-\left(-\dfrac{1}{2}\right)}
    \end{equation*}
    显然\(\lim\limits_{n \to\infty} S(n) = \dfrac 2 3\), 故原级数收敛.
\end{solution}
\begin{subproblem}
    $\sum\limits^\infty_{n = 1} \dfrac{1}{\sqrt[n]{n}}$
\end{subproblem}
\begin{solution}
    
    \begin{equation*}
        \because\forall n > 1, \dfrac 1 {\sqrt[n]{n}} > \dfrac 1 n,\ 
        \therefore \sum\limits^\infty_{n = 1} \dfrac{1}{\sqrt[n]{n}}
            > \sum\limits^\infty_{n = 1} \dfrac{1}{n}
    \end{equation*}
    而调和级数是发散的,\ 所以根据比较准则,\ 原级数发散.
\end{solution}
\begin{subproblem}
    \begin{equation*}
        \sum\limits^\infty_{n=1} \left(\dfrac 1 n + \dfrac 1 {2^n}\right)
    \end{equation*}
\end{subproblem}
\begin{solution}
    \begin{equation*}
        \sum\limits^\infty_{n=1} \left(\dfrac 1 n + \dfrac 1 {2^n}\right) 
        = \sum\limits^\infty_{n=1} \dfrac 1 n + \sum\limits^\infty_{n=1} \dfrac 1 {2^n}
    \end{equation*}
    由于调和级数是发散的,\ 等比级数是收敛的,\ 故原级数发散.
\end{solution}
\begin{problem}
    求八进制无限循环小数$\left(36.07360736\cdots\right)_8$的值.
\end{problem}
\begin{solution}
    原数等于:
    \begin{align*}
        & \left(3 \times 8^3 + 6 \times 8^2 + 7 \times 8^0\right) \times
        \sum\limits^\infty_{n=0}\dfrac 1 {8^{4n+2}}\\
        = & (1536+384+7) \times \dfrac{1}{64}\lim\limits_{n \to \infty}
        \cdot\dfrac{1 - \dfrac{1}{4096^n}}{1-\dfrac{1}{4096}}\\
        = & \dfrac {123328} {4095}
    \end{align*}
\end{solution}
\begin{problem}
    证明:\ 级数$\sum\limits^\infty_{n=1}\left(u_{n+1}-u_n\right)$收敛的充分必要条件是
    $\{u_n\}$收敛.
\end{problem}
\begin{solution}
    这个级数的部分和为
    \begin{align*}
        S(n) = \sum\limits^\infty_{n=1}\left(u_{n+1}-u_n\right) = u_n - u_1\\
    \end{align*}
    故
    \begin{align*}
        &\lim\limits_{n \to \infty} \sum\limits^\infty_{n=1}\left(u_{n+1}-u_n\right)
        = \lim\limits_{n \to \infty} u_n - u_1\\
        \Rightarrow & \lim\limits_{n \to \infty} u_n 
        = \sum\limits^\infty_{n=1}\left(u_{n+1}-u_n\right) + u_1
    \end{align*}
    由于$u_1 \ne \infty$, 故当$\{u_n\}$收敛时,\ 级数也收敛,\ 反之亦然.
\end{solution}
\begin{problem}
    用比较法确定以下级数的敛散性:
\end{problem}
\begin{subproblem}
    \(\sum\limits^\infty_{n=2} \dfrac{1}{(n-1)(n+2)}\)
\end{subproblem}
\begin{solution}
    由于$\forall n >= 2, \dfrac{1}{(n-1)(n+2)} < \dfrac{1}{n(n-1)} 
    < \left(\dfrac{1}{n-1} - \dfrac{1}{n}\right)$,\ 所以:
    \begin{align*}
        & \sum\limits^\infty_{n=2} \dfrac{1}{(n-1)(n+2)}
        < \sum\limits^\infty_{n=2} \dfrac{1}{n(n-1)}
        < \sum\limits^\infty_{n=2} \left(\dfrac{1}{n-1} - \dfrac{1}{n}\right)
        = \lim\limits_{n \to \infty} 1-\dfrac 1 n = 1\\
    \end{align*}
    所以原级数收敛.
\end{solution}
\begin{subproblem}
    $\sum\limits^\infty_{n=1} \tan \dfrac{\pi}{4n}$
\end{subproblem}
\begin{solution}
    易知$\forall x \in \left( 0, 2\pi \right), \tan x > x$, 所以:
    \begin{equation*}
        \sum\limits^\infty_{n=1} \tan \dfrac{\pi}{4n}
        > \sum\limits^\infty_{n=1} \dfrac{\pi}{4n}
        > \dfrac{\pi}{4}\sum\limits^\infty_{n=1} \dfrac 1 n
    \end{equation*}
    由于调和级数是发散的,\ 所以根据比较准则,\ 原级数发散.
\end{solution}
\begin{subproblem} 
    $\sum\limits^\infty_{n=1}\dfrac{1}{n\sqrt[n]{n}}$
\end{subproblem}
\begin{solution}
    $\forall n > 2, n\sqrt[n]{n} < 2n 
    \Rightarrow \forall n>2, \dfrac{1}{n\sqrt[n]{n}} > \dfrac{1}{2n}$, 所以:
    \begin{equation*}
        \sum\limits^\infty_{n=1}\dfrac{1}{n\sqrt[n]{n}} 
        < \sum\limits^\infty_{n=1}\dfrac{1}{2n}
    \end{equation*}
    由于调和级数是发散的,\ 所以根据比较准则, \ 原级数发散.
\end{solution}
\begin{subproblem}
    $\sum\limits^\infty_{n=1} 
    \left( 2n - \sqrt{n^2+1} - \sqrt{n^2-1} \right)$
\end{subproblem}
\begin{solution}
    注意到\(\forall n>=1, \sqrt{n^2-1} < n\), 故
    \begin{align*}
        2n - \sqrt{n^2+1} - \sqrt{n^2-1} & > n - \sqrt{n^2-1} 
        = \dfrac{1}{n+\sqrt{n^2+1}} > \dfrac{1}{2n}
    \end{align*}
    所以\(\sum\limits_{n=1}^\infty \left(2n - \sqrt{n^2+1} - \sqrt{n^2-1}\right)
    > \sum\limits_{n=1}^\infty \dfrac{1}{2n}\).
    由于调和级数是发散的, 所以根据比较准则, 原级数发散.
\end{solution}
\begin{problem}
    用检比法判断以下级数的敛散性:
\end{problem}
\begin{subproblem}
    \(\sum\limits^\infty_{n=1}\dfrac{n^2}{3^n}\)
\end{subproblem}
\begin{solution}
    令$a_n = \dfrac{n^2}{3^n}$,\ 则
    \[
        \lim\limits_{n\to\infty}\dfrac{a_{n+1}}{a_n} 
        = \lim\limits_{n\to\infty}\dfrac{(n+1)^2}{3n^2}
        = \dfrac{1}{3}
    \]
    由检比法可知改级数收敛.
\end{solution}
\begin{subproblem}
    \(\sum\limits^\infty_{n=1} n\tan\dfrac{\pi}{2^{n+1}}\)
\end{subproblem}
\begin{solution}
    令\(a_n = n\tan\dfrac{\pi}{2^{n+1}}\), 则
    \[
        \lim\limits_{n\to\infty}\dfrac{a_{n+1}}{a_n}
        = \lim\limits_{n\to\infty}
            \dfrac{(n+1)(1-\tan\dfrac{\pi}{2^{n+2}})}{2n}
        = \dfrac 1 2 
    \]
    由检比法可知原级数收敛.
\end{solution}
\begin{subproblem}
    \(
        \sum\limits^\infty_{n=1} 
        \dfrac{a^n}{(1+a)(1+a^2)\cdots(1+a^n)}
        (a>0)
    \)
\end{subproblem}
\begin{solution}
    令$t_n=\dfrac{a^n}{(1+a)(1+a^2)\cdots(1+a^n)}$,则
    \[
        \lim\limits_{n\to\infty}\dfrac{t_{n+1}}{t_n}
        = \lim\limits_{n\to\infty}\dfrac{a}{1+a^n}
    \]
    当$a<1$时,上式极限为$a$;当$a\ge 1$时,上式极限为0.\par
    由检比法可知原级数收敛.
\end{solution}
\end{document}